\documentclass{article}
\usepackage[margin=2cm]{geometry}
\usepackage{amsmath}
\usepackage{graphicx}
\usepackage{booktabs}
\usepackage[utf8]{inputenc}
\usepackage{ragged2e}
\setlength{\emergencystretch}{3em}
\begin{document}
\sloppy

where NAB, SAB, and MAB are again given by, but now with the appropriate redefinition
of the coset representative as per. It should be noted that while the FDA analysis presented.
in Appendix is a strong motivation for the form of the supersymmetry variations presented
above, it is not a proof. To actually derive the form of these variations, one must first
introduce curvature terms representing deviations from zero of each line in the free differential
algebra. An application of the exterior derivative to the resulting expressions then gives rise
to Bianchi identities, which must be solved before obtaining the explicit form of the fermion
variations. This is a rather involved process, and so for the moment we will content ourselves
with the motivating comments provided by the FDA. We will take the eventual presence of
smooth supersymmetric solutions consistent with the equations of motion as a posteriori.
evidence for the legitimacy of these variations. A nice property of the variations above is
the fact that they are consistent with the following SO(6)-invariant symplectic Majorana
condition,


\begin{equation}
\bar { \psi } _ { A } = \epsilon ^ { A B } \psi _ { B } ^ { T } \mathcal { C } \qquad ( \mathbf { l } )
\end{equation}


The consistency of such a condition allows us to work with symplectic Majorana spinors just
as in the Lorentzian case, though the symplectic Majorana condition utilized here is different
than that of the Lorentzian case. As mentioned before, we will be concerned with only the
simplest case of a single non-zero SU(2)r-charged vector multiplet scalar \$3, i.e. we take
\$1 = \$2 = 0. It can be easily verified that this is a consistent truncation, and is in fact the
most general choice of non-vanishing fields that can preserve SO(4,2)  U(1)R. With this
consistent truncation, the functions NAB, SAB, and MAB appearing in the supersymmetry
variations reduce to


\begin{equation}
\begin{array} { c } { S _ { A B } = i S _ { 0 } \epsilon _ { A B } + i S _ { 3 } \gamma ^ { 7 } \sigma _ { A B } ^ { 3 } } \\ { N _ { A B } = - N _ { 0 } \epsilon _ { A B } - N _ { 3 } \gamma ^ { 7 } \sigma _ { A B } ^ { 3 } } \\ { M _ { A B } ^ { I } = M _ { 9 } \gamma ^ { 7 } \epsilon _ { A B } + M _ { 3 } \sigma _ { A B } ^ { 3 } \qquad ( 2 ) } \end{array}
\end{equation}


where we have defined


\begin{equation}
\begin{array} { c } { S _ { 0 } = \frac { 1 } { 4 } \left ( g \cos \phi ^ { 3 } e ^ { \sigma } + m e ^ { - 3 \sigma } \cosh \phi ^ { 0 } \right ) } \\ { \qquad }\end{array}
\end{equation}


Importantly, note that S3, N3, and M3 are now purely imaginary, in contrast to the Lorentzian.
case . In all that follows we will set m = 1/2n such that the radius of AdS6 is one.

\end{document}
